% SHARD-ID: AQM-51-REGULAR-FUNCTION-AFFINE-SYMMETRY-ACTION
% SHARD-TITLE: Affine Symmetry Action on Regular Functions
% SHARD-SUMMARY: Defines the pushforward action \(g_\#F=F\circ g^{-1}\) on regular functions over standard opens.
% SHARD-SUMMARY: Proves compatibility with residue profiles, finite Weyl labels, and reduced smeared operators.
% SHARD-SUMMARY: Separates the canonical Artin-ring transport from any extra thickened Weyl character convention.
% SHARD-KEYWORDS: regular function, affine symmetry, residue profile, Weyl net, smearing, Artin ring

\section{Affine Symmetry Action on Regular Functions}
\label{sec:regular-function-affine-symmetry-action}

\subsection{Status and Source Anchors}
This shard extends AQM-46 and AQM-49.  The group
\(\operatorname{Aut}_k(\mathbf A^1_k)\), its action on closed points, and the
residue-field isomorphisms \(\tau_{g,\pi}\) are those of AQM-46.  Regular
functions, residue profiles, reduced supports, and reduced smearing are those
of AQM-49.  The Artin-ring layer is AQM-36.  No new external source is used;
the new assertions below are checked algebraic consequences of those shards.

\subsection{Lamport Dependency Skeleton}
\[
\begin{array}{rcl}
D1&:&k=\mathbb F_q,\quad R=k[x],\quad X=\Spec R,\quad
      g=g_{a,b}:x\mapsto ax+b.\\
D2&:&\theta_g(f)(x)=f(ax+b),\quad
      \sigma_g=\theta_{g^{-1}},\quad F^g=\sigma_g(F).\\
L1&:&\sigma_g:R_h\to R_{\sigma_g(h)}\text{ is an isomorphism and }
      gD(h)=D(\sigma_g(h)).\\
D3&:&\operatorname{Prof}(U)=\prod_{x_\pi\in U}K_\pi^2.\\
D4&:&(S_g^\Pi e)_{\pi^g}=
      (\tau_{g,\pi}^{-1}q_\pi,\tau_{g,\pi}^{-1}p_\pi).\\
T1&:&\operatorname{ev}_{gU}(F^g,G^g)=
      S_g^\Pi\operatorname{ev}_U(F,G).\\
T2&:&\text{The affine action preserves the distinction between profiles and
      finite Weyl labels.}\\
D5&:&h^g=\sigma_g(h),\quad
      W_h^{\rm red}(F,G)=W_{Z_{\rm cl}(h)}(e_h^{\rm red}(F,G)).\\
T3&:&\alpha_g(W_h^{\rm red}(F,G))=
      W_{h^g}^{\rm red}(F^g,G^g).\\
D6&:&A_{\pi,e}=R/(\pi^e),\quad
      A_{\pi^g,e}=R/((\pi^g)^e).\\
T4&:&\sigma_g\text{ induces }A_{\pi,e}\simeq A_{\pi^g,e}
      \text{ and preserves the multiplicity }e.
\end{array}
\]

\subsection{Pushforward of Regular Functions D1--L1}
Let \(g=g_{a,b}\), with \(a\in k^\times\) and \(b\in k\).  Its pullback on the
coordinate ring is
\[
  \theta_g:R\longrightarrow R,\qquad \theta_g(f)(x)=f(ax+b).
\]
For regular functions we use the covariant, or pushforward, notation
\[
  F^g=g_\#F:=\sigma_g(F),\qquad \sigma_g=\theta_{g^{-1}}.
\]
Thus, as a function on points,
\[
  (g_\#F)(y)=F(g^{-1}y).
\]

\paragraph{Lemma \(L1\).}
If \(h\ne0\), then \(\sigma_g\) extends to a ring isomorphism
\[
  \sigma_g:R_h\longrightarrow R_{\sigma_g(h)}
\]
by
\[
  \sigma_g(u/h^m)=\sigma_g(u)/\sigma_g(h)^m.
\]
Moreover
\[
  g(D(h))=D(\sigma_g(h)).
\]

\emph{Proof.}
The map \(\sigma_g\) is a \(k\)-algebra automorphism of \(R\), with inverse
\(\theta_g\).  It sends the multiplicative set
\(\{1,h,h^2,\ldots\}\) bijectively to
\(\{1,\sigma_g(h),\sigma_g(h)^2,\ldots\}\), so it extends uniquely to the
displayed localization isomorphism.

For a point \(y\in X\),
\[
  y\in g(D(h))
  \Longleftrightarrow
  g^{-1}y\in D(h)
  \Longleftrightarrow
  h(g^{-1}y)\ne0
  \Longleftrightarrow
  \sigma_g(h)(y)\ne0.
\]
This is exactly \(y\in D(\sigma_g(h))\).

Composition has the expected order:
\[
  (gh)_\#=g_\#h_\#,
\]
because \((gh)^{-1}=h^{-1}g^{-1}\).

\subsection{Residue Profiles D3--T2}
For an open \(U\), define the full residue-profile set
\[
  \operatorname{Prof}(U)=\prod_{x_\pi\in U}K_\pi^2.
\]
The finite-support Weyl label group \(E(U)\) of AQM-46 and AQM-49 is the
subgroup of \(\operatorname{Prof}(U)\) with finite closed support.

Define
\[
  S_g^\Pi:\operatorname{Prof}(U)\longrightarrow\operatorname{Prof}(gU)
\]
by
\[
  (S_g^\Pi e)_{\pi^g}
  =
  \bigl(\tau_{g,\pi}^{-1}(q_\pi),
        \tau_{g,\pi}^{-1}(p_\pi)\bigr),
\]
where \(e_\pi=(q_\pi,p_\pi)\).  On finite-support profiles this is the
finite Weyl-label map \(S_g\) from AQM-46.

\paragraph{Theorem \(T1\).}
For \(F,G\in\Gamma(U,\mathcal O_X)\),
\[
  \operatorname{ev}_{gU}(F^g,G^g)=
  S_g^\Pi\operatorname{ev}_U(F,G).
\]

\emph{Proof.}
It is enough to check one closed point \(x_\pi\in U\).  Let
\(x_{\pi^g}=g(x_\pi)\).  The residue map
\(\tau_{g,\pi}:K_{\pi^g}\to K_\pi\) is induced by \(\theta_g\).  Therefore,
for every \(H\in\Gamma(gU,\mathcal O_X)\),
\[
  \tau_{g,\pi}(H(x_{\pi^g}))=(\theta_gH)(x_\pi).
\]
Apply this to \(H=F^g=\sigma_g(F)\).  Since \(\theta_g\sigma_g=\operatorname{id}\),
\[
  \tau_{g,\pi}(F^g(x_{\pi^g}))=F(x_\pi).
\]
Thus \(F^g(x_{\pi^g})=\tau_{g,\pi}^{-1}(F(x_\pi))\).  The same argument applies
to \(G\), proving the displayed profile identity.

\paragraph{Theorem \(T2\).}
The affine action never turns a nonzero regular-function profile into a
finite Weyl label, and never turns the zero profile into a nonzero one.

\emph{Proof.}
The map \(S_g^\Pi\) is a bijection of closed-point profiles.  It sends the zero
profile to the zero profile.  It sends a finite-support profile to a
finite-support profile because \(x_\pi\mapsto x_{\pi^g}\) is a bijection of
closed points.  It sends an infinite-support profile to an infinite-support
profile for the same reason.  AQM-49 proves that a nonzero regular-function
pair on a nonempty standard open has cofinite, hence infinite, closed support.
Therefore the profile/operator distinction is invariant under \(g\).

\subsection{Reduced Smeared Operators D5--T3}
For a nonzero polynomial \(h\), put
\[
  h^g=\sigma_g(h).
\]
This polynomial need not be monic.  Multiplying it by a nonzero scalar does
not change \(D(h^g)\), \(Z_{\rm cl}(h^g)\), or the reduced Weyl support.  If
\(\pi\) is monic irreducible, then \(\pi\mid h\) exactly when
\(\pi^g\mid h^g\), so
\[
  g(Z_{\rm cl}(h))=Z_{\rm cl}(h^g).
\]

\paragraph{Theorem \(T3\).}
For \(F,G\in R\),
\[
  \alpha_g\!\left(W_h^{\rm red}(F,G)\right)
  =
  W_{h^g}^{\rm red}(F^g,G^g).
\]

\emph{Proof.}
At the source point \(x_\pi\in Z_{\rm cl}(h)\), the reduced label is
\[
  (F\bmod\pi,\ G\bmod\pi)\in K_\pi^2.
\]
At the target point \(x_{\pi^g}\), Theorem \(T1\) gives the label
\[
  (F^g\bmod\pi^g,\ G^g\bmod\pi^g)
  =
  \bigl(\tau_{g,\pi}^{-1}(F\bmod\pi),
        \tau_{g,\pi}^{-1}(G\bmod\pi)\bigr).
\]
Thus
\[
  e_{h^g}^{\rm red}(F^g,G^g)=S_g e_h^{\rm red}(F,G).
\]
AQM-46 proves
\[
  \alpha_g(W(e))=W(S_ge)
\]
on finite supports.  Substituting the displayed labels proves the theorem.

Consequently the rule is simple: move the support by \(g\), and precompose the
regular functions by \(g^{-1}\).

\subsection{Artin-Ring Transport D6--T4}
Let \(h=\prod_i\pi_i^{e_i}\).  AQM-49 records that the reduced theory sees
only the distinct \(\pi_i\), while the thickened layer uses
\[
  A_{\pi_i,e_i}=R/(\pi_i^{e_i}).
\]

\paragraph{Theorem \(T4\).}
For each \((\pi,e)\), the automorphism \(\sigma_g\) induces a ring
isomorphism
\[
  \bar\sigma_{g,\pi,e}:R/(\pi^e)\longrightarrow R/((\pi^g)^e).
\]
It sends the class of \(F\) to the class of \(F^g\), and it preserves the
exponent \(e\).

\emph{Proof.}
AQM-46 gives
\[
  \pi^g=a^{\deg\pi}\sigma_g(\pi).
\]
Hence the ideals \((\sigma_g(\pi)^e)\) and \(((\pi^g)^e)\) are equal, because
they differ by a unit.  Therefore \(\sigma_g\) descends to the displayed
quotient isomorphism.  The exponent remains \(e\) because an isomorphism sends
the \(e\)-th power of the source ideal to the \(e\)-th power of the target
ideal.

This is the canonical thickened support transport.  A thickened Weyl
automorphism also needs a character statement.  If the target generating
character is transported through \(\bar\sigma_{g,\pi,e}\), the same proof as
AQM-36 gives Weyl covariance.  If one insists on separately chosen default
top-coefficient characters at source and target, a dual momentum-label
correction must be supplied before claiming character preservation.

\subsection{Conclusion}
The affine symmetry action on regular functions is:
\[
  (h,F,G)\longmapsto(h^g,F^g,G^g)
  =
  (\sigma_g(h),\sigma_g(F),\sigma_g(G)).
\]
On profiles it is the product residue relabelling \(S_g^\Pi\).  On finite
reduced Weyl smearings it is exactly the local-net automorphism of AQM-46:
\[
  \alpha_gW_h^{\rm red}(F,G)=W_{h^g}^{\rm red}(F^g,G^g).
\]
Thus the symmetry does not add compatibility conditions between closed-point
qudits.  It only permutes the closed supports and relabels each residue field
by the induced field isomorphism.
